%\input ../util/bookmacros
%\input ../util/docparams

\appendix{Appendix D}{Mathematical Induction}

\vskip -.075in
\section{Domino Theory}
\vskip -.1in
 Several times in the main text we claimed a formula for {\it every\/} $n$
-- for example $(x^n)^\prime = n\, x^{n-1}$, or $\sum_{i=1}^n i = n(n+1)/2$ --
after checking only the first few cases. How can a finite amount of work
prove infinitely many statements? That is exactly what Mathematical
Induction buys us.
Mathematical Induction is a technique used to determine the truth of a sequence of statements.
The sequence is usually
large or infinite. Normally, determining the truth of such a collection would be
time consuming -- if not impossible. However,
if there is a strong connection between successive elements of the sequence, then it can be much easier.
By connection I mean that when a given statement in the sequence 
is true it is relatively easy to show that the next statement is true.

In this sense, Induction is a generalization of what might be called ``The Domino Theory''. 
If you have 
ever played with dominoes then, in a sense, you already understand Mathematical 
Induction. To see why, consider placing a number of (upright) dominoes in a room. Knocking 
them down may be easy or hard depending on how they are laid out. If they 
are put next to each other in the ``standard domino way'', then knocking 
one domino down will knock all the rest down. How would you describe this to 
a person who has never seen dominoes fall in this way? You might start by 
labeling the dominoes; that is, after setting up the dominoes they have a natural ordering -- they form a sequence. 
Let's label them $D_1, D_2, \ldots, D_n$. You then 
tell the domino newcomer the dominoes are arranged so that whenever any one 
domino falls the domino after it falls. By this we mean domino $D_2$ is 
placed after domino $D_1$ so that when $D_1$ falls $D_2$ falls. In general, 
if domino $D_k$ falls then $D_{k+1}$ falls as well, if $1\le k< n$. 
Now you ask the newcomer: 
``Without looking into the room with the dominoes, if I told you that the 
first domino, $D_1$, has fallen, what can you say about the rest of dominoes?''
So, you've given the newcomer two facts:
\beginEnum
\enum{Domino $D_1$ has fallen.}
\enum{Whenever domino $D_k$ falls, domino $D_{k+1}$ falls 
provided $1 \le k< n$. That is, when one domino falls the next in line falls.}
\endEnum
This is all he needs to conclude that all of the dominoes in the room have fallen.

This is the essence of mathematical Induction. Mathematical Induction concerns the 
{\it truth\/} of propositions labeled $P_1, P_2, \ldots, P_n$. 
Like Dominoes, the propositions are either easy or hard to prove depending 
on how they are ``arranged''.
If the propositions 
are arranged like dominoes; that is, if
\beginEnum
\enum{Proposition $P_1$ is true;}
\enum{Whenever proposition $P_k$ is true, proposition $P_{k+1}$ is true
provided $1 \le k< n$;}
\endEnum
then the conclusion is that all of the propositions are true. 

In moving 
from dominoes to propositions 
all we've done is to replace statements like 
``domino $D_k$ falls'' with ``proposition $P_k$ is true''.

This result extends to the case where there are an infinite number of 
propositions (dominoes).
In this case the induction conditions (domino conditions) are:
\beginEnum
\enum{Proposition $P_1$ is true.}
\enum{Whenever proposition $P_k$ is true, proposition $P_{k+1}$ is true
provided $1 \le k$.}
\endEnum
The conclusion under these conditions is that {\it all\/} of the $P$ 
statements are true.

One may ask how easy is it to show that any given proposition implies the next? Just as one may ask 
how one verifies that a given domino falling knocks the next one down?
For many collections of statements, $P_n$, it is the case that the proof of one proposition implying the next 
is generic, and essentially one only has to do one proof!

\section{Applications}
We apply Mathematical Induction to a few problems below. It is worth keeping 
in mind what must be shown in the second step of induction. 
For each $k \ge 1$ you must show: {\it if\/} $P_k$ is true,
{\it then\/} $P_{k+1}$ is true. You are not assuming that $P_k$ actually is
true -- this step alone does not tell us that $P_k$ is
true or that $P_{k+1}$ is true. If you get confused think of dominoes.
To show domino 
$D_{k+1}$ falls when $D_{k}$ falls you would need to {\it assume\/}  
domino $D_k$ has fallen and show that $D_{k+1}$ falls as a consequence.

\subsection{A Formula for Adding the First $n$ Integers}
By adding up the numbers from 1 to 2, then 1 to 3, and so forth, someone guesses 
that the formula for the sum of the first $n$ integers is:
$$
\sum_{i=1}^n i = n(n+1)/2
$$
How can one prove this? We do it by Mathematical Induction. Let $P_n$ be the statement%
\footnote{\kern 0.5pt \raise 0.5ex \hbox{*}}{We surround statements in double quotes to indicate that they are 
statements with no assigned truth value. You can think of our usual mathematical statements such as, $3 = 2 - 1$, as consisting
of two things: The first is the statement ``3 is equal to 2 minus 1''. The second thing that is conveyed is that the 
statement is true. In this section we wish to make the separation of a statement from its truth value. So we will refer to 
statements, $P_n$, in quotes to indicate that they are statements without (yet) an assigned truth value.}
``$\sum_{i=1}^n i = n(n+1)/2$''. We want to show that every statement, $P_1, P_2, \ldots, $ is true. 
Mathematical Induction tells us that if $P_1$ is true 
and it is the case that whenever $P_k$ is true it follows that $P_{k+1}$ is 
true as well, then all of the statements in the sequence $\{P_n\}_{n=1}^\infty$ 
are true.

First we show $P_1$ is true. But $P_1$ is the statement 
``$\sum_{i=1}^1 i = 1(1+1)/2$''. The left hand side is 1 and the right hand side 
is 1; therefore, $P_1$ is true, since ``1 = 1'' is a true statement.

The next thing to do is to {\it assume\/} for any given $k$ ($k \ge 1$) $P_k$ is true. This 
means we assume ``$\sum_{i=1}^k i = k(k+1)/2$'' is a true statement; that is we assume: $\sum_{i=1}^k i = k(k+1)/2$.
We then use this to show 
``$\sum_{i=1}^{k+1} i = (k+1)(k+2)/2$'' is a true statement; that is, $\sum_{i=1}^{k+1} i = (k+1)(k+2)/2$. 
We proceed as follows: Take the sum
on the left and relate it to $P_k$.
$$
\sum_{i=1}^{k+1} i = (k+1) + \sum_{i=1}^{k} i 
$$
We use the assumed truth of $P_k$ to rewrite the sum on the right hand side so 
that 
$$
\sum_{i=1}^{k+1} i = (k+1) + \sum_{i=1}^{k} i = (k+1) + k(k+1)/2 = 
\bigl(2(k+1) + k(k+1) \bigr)/2 = (k+1)(k+2)/2
$$
That is,
$$
\sum_{i=1}^{k+1} i = (k+1)(k+2)/2
$$
We have just shown the truth of the statement  $P_{k+1}$,  ``$\sum_{i=1}^{k+1} i = (k+1)(k+2)/2$'', 
from the assumed truth of $P_k$. We see that
we have done so for every $k \ge 1$; the same argument works for {\it every\/} $k$.%
\footnote{\kern 0.5pt \raise 0.5ex \hbox{\dag}}{Notice that this is the same {\it magic\/} we have when finding formulas for the derivative of
certain functions; namely, we find a way to compute the formula for the derivative at a given ``x'' and that same formula works 
for all other ``x''s. If this wasn't the case for a lot of elementary functions, calculus wouldn't be useful. This is the same for induction:
If the way to show the truth of statement $P_k$ implies that $P_{k+1}$ is true {\it changes} for each $k$, then induction 
isn't useful. This is analogous to writing a function in a programming language that is parameterized. Sometimes rather than writing
many such functions, one can write {\it one function\/} that takes one or more parameters while the body of
the function is {\it generic\/}. This generic logic doesn't
always happen in programming nor does it always happen in mathematics, but when it does and we can identify the generic argument,
Induction gives us tremendous power.}
Mathematical Induction tells us that $P_n$ is true for all $n$. That is, 
we have shown that 
$$
\sum_{i=1}^n i = n(n+1)/2 \quad {\rm for} \; 1 \le n
$$
Notice, by the way, that the inductive step amounts to checking that
$\Delta$ of both sides agree -- induction is the engine that makes our
telescoping (discrete Fundamental Theorem) arguments rigorous.

Notice what helped here was that we had a natural ordering of the propositions
(like our room of dominoes) and
showing the truth of a given $P_k$ implying $P_{k+1}$ was the same for all $k$. Most successful 
applications of induction behave similarly.
If the proofs were different for each $k$, then
induction wouldn't be as useful; we would have a lot of work to do.

\subsection{The Derivative of the Power Functions}
We know the derivative of the function $x$; it is $1$. 
What is the derivative of $x^n$ when $n$ is an integer? We stated earlier in the book 
that the answer is $(x^n)^\prime = n\, x^{n-1}$. We can use Mathematical 
Induction to show this for $n\ge1$. Let $P_n$ be the statement: 
``$(x^n)^\prime = n \, x^{n-1}$'' First, 
$(x^1)^\prime = x^\prime = 1$; while $1 x^{1-1} = 1 * x^0 = 1 * 1 = 1$. 
So, $P_1$ is true.
We {\it assume\/} now for any $k$ ($k\ge1$), $P_k$ is true and try to show $P_{k+1}$ is true.
Statement $P_k$ true means $(x^k)^\prime = k \, x^{k-1}$. Now, 
the derivative of $x^{k+1}$ can be related to $x^k$ since
$x^{k+1} = x \, x^k$. Using the product rule and the assumed truth of the formula for the derivative 
of $x^k$ gives:
$(x x^k)^\prime = x^\prime x^k + x (x^k)^\prime = 1 x^k + x (k\, x^{k-1}) 
= x^k + k\,x^k = (k+1)x^k$. Or, $(x^{k+1})^\prime = (k+1)x^{(k+1)-1}$. 
That is, the statement $P_{k+1}$ 
is true when $P_k$ is true (for {\it any\/} $k\ge 1$). By Mathematical Induction, $P_n$ is true 
for all $n\ge 1$. That is,
$$
(x^n)^\prime = n\, x^{n-1} \quad {\rm for} \; n\ge 1
$$
What about when $n$ is $0$ or negative? We haven't proved the formula in those cases. One can easily check the formula for
$n=0$. For negative $n$ we can use induction again. In this case let $P_n$ be the statement ``$(x^{-n})^\prime = -n x^{-n-1}$'' for
$n \ge 1$ and $x\ne0$. Induction requires statements labeled from 1, so we index the statements by $n \ge 1$, where $P_n$ concerns the power $-n$.
To check $P_1$ -- the claim $(x^{-1})^\prime = -x^{-2}$ -- one can use the definition of the derivative directly.
This requires work, but once done, one can mimic the induction proof above. Given that $P_1$ has been proved true we next show
that $P_k$ implies $P_{k+1}$. Consider an $x \ne 0$, then 
$$
(x^{-(k+1)})^\prime = (x^{-1} x^{-k})^\prime = (x^{-1})^\prime x^{-k}
+ x^{-1} (x^{-k})^\prime
$$ 
The last statement is true by the product rule. Using the assumed truth of $P_1$ and $P_k$ we may write
$$
(x^{-(k+1)})^\prime = -1x^{-1-1} x^{-k}
+ x^{-1} (-kx^{-k-1}) = -x^{-k-2} - kx^{-k-2} = -(k+1) x^{-(k+1) - 1}
$$
We have just shown that $P_{k+1}$ is true given that $P_k$ is true. By mathematical induction $P_n$ must be true for $n\ge 1$.
The only thing we left out was the proof that $P_1$ is true. We leave this as an exercise for the reader.

